\section{Reflection and rotation groups}\label{sec:groups}
For $n\ge3$ write $G=D_n$ for the symmetry group of a regular
$n$-gon, of order $2n$, and $V$ for its real representation on the
plane. Fix a connected principal unramified $G$-cover $p:Z\to C$.
It determines the flat plane bundle $\mathcal V=Z\times_G V$.
The twist $D^{\mathcal V}$ acts in a flat real frame by $D$ on each
component. Equivalently its kernel is
$\operatorname{Hom}_G(V,\ker p^*D)$, using the invariant Euclidean
metric on $V$. Its real Fredholm index is zero. Choose one subgroup
$H_i$ generated by a reflection in each conjugacy class of such
subgroups. Set
\[
    f_i:Y_i=Z/H_i\longrightarrow C,\qquad
    B_{i,t}=f_i^*D_t,\qquad A_t=p^*D_t.
\]
For odd $n$ there is one quotient $Y_i$; for even $n$ there are two.
Each has degree $n$ over $C$.

We say a path crosses the \emph{wall for $\mathcal V$} transversely
at zero if $\ker D_0^{\mathcal V}$ is a real line and the derivative
along the path induces an isomorphism from this kernel to its
cokernel. Generic means residual in the smooth topology on paths of
fixed-symbol operators. We consider crossings in the interior of
paths of base-invertible operators; endpoints of such paths can be
held fixed. In the following theorem a term whose indicated degree
is not a positive integer is taken to be zero.

\begin{theorem}\label{thm:dihedral}
At a transverse crossing of the wall for $\mathcal V$ on a generic
path, for every finite $M$ there is $\delta_M>0$ such that $D_t$ is
rigid through $M$ on $0<|t|\le\delta_M$. For
$-\delta_M\le t_-<0<t_+\le\delta_M$, $1\le d\le M$ and every $k$,
one has
\begin{equation}\label{eq:main}
  L^C_{d,k}(D_{t_+})-L^C_{d,k}(D_{t_-})
  =\sum_i \frac{2}{m_i} L^{Y_i}_{d/n,k}(B_{i,t_+})
           -L^Z_{d/(2n),k}(A_{t_+}), \qquad
              m_i=|N_G(H_i)/H_i|.
\end{equation}
Each $m_i$ equals one for odd $n$ and two for even $n$.
For every $p:Z\to C$ as above and every $N$ of the specified degree
there is a path with such a crossing satisfying all the conclusions.
\end{theorem}


\begin{proof}
Theorem~\ref{thm:kernels} gives central kernels on all smooth covers
and matching on nodal domains. The proof of bounded-degree isolation
in Section~\ref{sec:family} uses only this classification, fiberwise
injectivity and transversality on the prescribed summand, so applies
to these paths too. The unit circle of the reflection plane has, as a
$G$-cell decomposition, two vertex orbits and one orbit of open edges,
obtained by subdividing along the reflection axes. Its class is
$\sum_i(2/m_i)[G/H_i]-[G/1]$. The proof of Theorem~\ref{thm:mixed}
now gives the asserted signed integrals, including the regular-cover
subtraction. In particular their balanced products obey
\begin{equation}\label{eq:prodwall}
 \frac{P_C^+(q)}{P_C^-(q)}
 =\frac{\prod_i P_{Y_i}^+(q^n)^{2/m_i}}{P_Z^+(q^{2n})}
       \pmod {q^{M+1}}
\end{equation}
by Corollary~\ref{cor:cutoff}.
\end{proof}
\begin{corollary}\label{cor:sign}
Suppose $n$ is odd and $M\ge2n$. Write $Y=Y_1$,
$B_t=f_1^*D_t$, $h=g(Y)$, $H_d=1+d(g-1)$ and
$\eta=\epsilon(B_{t_+})$. Let $\lambda_0\in H^1(Y;\mathbb F_2)$
classify the connected unramified double $Z\to Y$. Let $N_-$ count
the $\lambda\in H^1(Y;\mathbb F_2)\setminus\{0,\lambda_0\}$ for which
the twist $B_0^\lambda$ by the associated real sign line has negative
determinant sign. These twists are invertible. With $\Delta$ denoting
the value at $t_+$ minus the value at $t_-$, one has
\begin{equation}\label{eq:sign}
       \Delta n^C_{n,H_n}=2\eta,\qquad
       \Delta n^C_{2n,H_{2n}}=-\eta(1+2N_-).
\end{equation}
In every degree strictly below $n$ all exponent changes vanish.
\end{corollary}

\begin{proof}
From \eqref{eq:prodwall}, nonmultiples of $n$ have zero change.
The degree-one contribution on $Y$ at genus $h$ is $\eta$, giving
the first assertion. To compute the next multiple, all genus-$H_{2n}$
maps in degree two over $Y$ are connected unramified doubles. Indeed
Riemann--Hurwitz on the positive normalizations and the arithmetic
genus formula force one such component, with no self-nodes or
contracted decorations. Each double has weight $1/2$. The
degree-one contribution over $Z$ in this genus is $\epsilon(A_{t_+})$.
It cancels the contribution in twice the upstairs degree-two count
from the distinguished double $\lambda_0$.

Identify the reflection subgroup $H$ with $\mathbb F_2$. On
$\pi_1(Y)$ the monodromy of $Z\to Y$ is the character $\lambda_0$.
Thus the monodromy image in $G$ of the double of $Y$ specified by
$\lambda\ne0$ is conjugate to $\lambda_0(\ker\lambda)\le H$. It is
trivial precisely when $\ker\lambda\subseteq\ker\lambda_0$.
Both kernels have index two, so this happens precisely for
$\lambda=\lambda_0$. Every other double has reflection image.
The central kernel
on that double is a single line from $Y$ by
Theorem~\ref{thm:kernels}. The odd twist on the double is thus
injective on $Y$, and invertible by index zero. Its sign is constant
near the crossing. The full determinant sign at $t_+$ is
$\eta\,\epsilon(B_0^\lambda)$ by splitting under the deck
involution with its complex reference orientation. Formula
\eqref{eq:main} now gives
\[
      \Delta L^C_{2n,H_{2n}}
                   =\eta(2^{2h}-2-2N_-).
\]
In the product only the changing degree-$n$ factor can contribute
from a proper divisor. Only its genus-$h$ term enters at
$H_{2n}=2h-1$, since twice a class of genus $s>h$ first contributes
in genus at least $2s-1$. The exponent at genus $h$ changes by
$2\eta$. A balanced elementary genus-$h$ curve has connected
degree-two coefficient $(2^{2h}-1)/2$ in the
minimum genus $2h-1=H_{2n}$. This follows either by expanding its
partition formula \cite[Corollaries~7.2 and~7.3]{bp08} or by counting the
unramified doubles with complex orientation. Subtracting that
coefficient times $2\eta$ gives \eqref{eq:sign}. In using the
minimum genus here, the same lower bound holds for the unchanging
and changing endpoint exponents: connected degree-$d$ maps to
a genus-$g$ target have genus at least $H_d$, and the connected
elementary series on its $j$-fold multiple starts in genus at least
$1+j(s-1)$ for $s\ge2$. Induction on the label and then genus gives
the exponent bound (Proposition~\ref{prop:factors}).
\end{proof}

\begin{corollary}\label{cor:separation}
Fix $C$ and $N$ as above. For every odd $n\ge3$ there are two
operators on $N$, rigid on all smooth covers, whose geometric local
products coincide through degree $n-1$ and differ in degree $n$.
They can be chosen arbitrarily close to the same transverse crossing.
Their degree-$2n$ exponents at genus $H_{2n}$ also differ.
\end{corollary}

\begin{proof}
In the surface group presentation send two of the $a$-generators
to a rotation and a reflection generating $D_n$, all the
$b$-generators to the identity, and the remaining $a$-generators
to the identity. This defines a surjection and hence a connected
unramified $D_n$-cover. Apply the realization in
Theorem~\ref{thm:kernels} and take $M=2n$. Choose two arbitrarily
close regular endpoints as in Theorem~\ref{thm:dihedral}.

Rigidity through $M$ is open here. Otherwise normalized kernel
graphs for operators approaching a rigid endpoint would have a
bounded-area limit of degree at most $M$ meeting positive height,
as in Proposition~\ref{prop:interval}. Any positive-height
component of that limit contradicts rigidity. In a small convex
neighborhood of each endpoint the connected Euler numbers therefore
remain constant through $M$, in every genus. All-cover rigidity is
dense in the unrestricted operator space. One can see this directly
from Section~\ref{sec:kernels}: omit the prescribed columns $s_1,s_2$
from the Petri test and use unrestricted local zeroth-order
perturbations. For each irreducible finite system the remaining
projective incidence has index at most $-1$ after moving its effective
branch values. The same compact exhaustions and Baire argument exclude
them all, including the trivial system and the prescribed plane.
Choose such an operator in each endpoint neighborhood. The full local
products exist by the compatible finite-energy constructions at every
cutoff. Their truncations through $2n$ agree with those of the chosen
endpoints by the connected comparison and the uniqueness argument of
Theorem~\ref{thm:mixed}. Corollary~\ref{cor:sign} gives the claim.
\end{proof}

\subsection{Two mixed representations}
The sign character and standard real plane of \(S_3\) have zero mixed class. The sphere of the sign line is \(G/A_3\) and the plane contributes \(2[G/C_2]-[G/1]\). The first Cartesian product with \(G/C_2\) is regular and the product with \(G/1\) is twice regular. Thus both single-wall transfers can be present while their mixed difference vanishes in every degree.

For the order-eight group \(G=\langle r,s\mid r^4=s^2=1,\ srs=r^{-1}\rangle\) take its standard plane and the sign line with \(r\mapsto-1\), \(s\mapsto1\). The plane sphere contributes \([G/\langle s\rangle]+[G/\langle rs\rangle]-[G/1]\). The sign kernel contains the first reflection and misses the second. Multiplying gives \(2[G/\langle s\rangle]-[G/1]\). Both examples are realized on every genus-\(g\) curve as above: choose a surjection of the surface group through a free group on two generators. The associated rank-three real bundle pulls back from a graph. Adding its determinant line produces an oriented rank-four bundle pulled back from that graph, hence trivial. The oriented real bundle underlying \(N\) is also trivial on the surface since \(c_1(N)\) is even. This gives the required monomorphism.

\subsection{The tetrahedral, octahedral and icosahedral walls}

Take \(W=\mathbb R^3\) and the rotation group of a regular tetrahedron, cube or icosahedron. Each standard real representation is faithful irreducible of real type: two nonparallel rotation axes of rotations of order greater than two preclude an invariant real line, and irreducibility in odd real dimension forces real commutant. The complement of the associated rank-three bundle in \(N_{\mathbb R}\), when an embedding exists, is a real line. The representation is oriented, so this complement must be orientable, hence trivial. Consequently a monomorphism exists exactly when \(w_2(E)=0\). Indeed \(N_{\mathbb R}\) is oriented with \(w_2=c_1(N)\bmod 2=0\), and oriented bundles of rank at least three over a surface are classified by \(w_2\). This condition is attained on every base curve of genus at least two: take the surface group quotient onto a free group on two generators and follow it by a generating pair for the rotation group. Its flat bundle pulls back from a graph, on which \(H^2\) vanishes.

Radially project the barycentric subdivision of the boundary of the solid onto \(S^2\). There are three vertex orbits (face center, edge midpoint, original vertex), three open edge orbits and two open triangle orbits. Edge interiors and triangle interiors have trivial stabilizer in the rotation group. The vertex stabilizer orders are the three entries in the table. Thus the equivariant Euler characteristic has one term for each vertex orbit and minus one regular orbit. In the table \(P_j\) denotes the positive-side geometric product on the indicated degree-\(j\) quotient, evaluated with \(q\) replaced by \(q^j\).

{\def\LTcaptype{table} % do not increment counter
\begin{longtable}[]{@{}
  >{\raggedright\arraybackslash}p{(\linewidth - 4\tabcolsep) * \real{0.3333}}
  >{\raggedright\arraybackslash}p{(\linewidth - 4\tabcolsep) * \real{0.3333}}
  >{\raggedright\arraybackslash}p{(\linewidth - 4\tabcolsep) * \real{0.3333}}@{}}
\toprule\noalign{}
\begin{minipage}[b]{\linewidth}\raggedright
Rotation group
\end{minipage} & \begin{minipage}[b]{\linewidth}\raggedright
Stabilizer orders
\end{minipage} & \begin{minipage}[b]{\linewidth}\raggedright
Quotient of the two downstairs products
\end{minipage} \\
\midrule\noalign{}
\endhead
\bottomrule\noalign{}
\endlastfoot
\(A_4\) & \(2,3,3\) & \(P_6 P_4^2/P_{12}\) \\
\(S_4\) & \(2,3,4\) & \(P_{12}P_8P_6/P_{24}\) \\
\(A_5\) & \(2,3,5\) & \(P_{30}P_{20}P_{12}/P_{60}\) \\
\end{longtable}
}

For \(A_4\) the two order-three vertex stabilizers are conjugate; their quotient targets give the same operator up to pullback equivalence over \(C\). For \(S_4\) the order-two stabilizer is the edge-midpoint half turn; on the four body diagonals it acts as a transposition. It is distinct from the subgroup generated by the square of the rotation of order four. The table proves that at minimum genus the first exponent change has magnitude two in tetrahedral degree four, and magnitude one in octahedral degree six and icosahedral degree twelve. The change in the tetrahedral exponent at degree six and its minimum genus also has magnitude one: no positive multiple of the degree-four term enters degree six. These are changes of geometric elementary exponents through a finite cutoff. For endpoints rigid in every degree, fix a cutoff containing the labels in question. Finite-degree rigidity is open by the normalized graph-cycle argument. Unrestricted all-cover rigidity is residual by the same projective index argument with no prescribed solution. Perturb each end within its small convex rigidity neighborhood. As in Proposition~\ref{prop:factors}, the full products exist and their exponents through the chosen cutoff persist. 

These conclusions use fiberwise monomorphism. At a rank-\(r\) input the first failure of real injectivity in a rank-four output has real codimension \(5-r\). The usual marked evaluation argument on a two-dimensional base excludes these defects automatically for a generic transverse rank-two crossing. In rank three the defect has codimension two, and in rank four it has codimension one. The construction gives realized crossing families and \eqref{eq:mixed} throughout their stated domain. Noninjective evaluations and ramified principal covers allow different orbit graphs to meet and require a separate wall calculation.
